\documentclass[12pt]{article}
\usepackage{graphicx} % Required for inserting images
\usepackage[utf8]{inputenc} % Codificación de caracteres
\usepackage[spanish]{babel} % Idioma del documento
\usepackage{amsmath, amssymb, amsthm}
\usepackage{pgfplots}
\usepackage{tabularx}
\usepackage[margin=1in]{geometry}
\usepackage{svg}
\usepackage{float}
\usepackage{fancyhdr}
\usepackage{setspace}
\usepackage[most]{tcolorbox}


\newtcolorbox{comentarioprofe}[1]{
    colback=blue!5,
    colframe=blue!50!black,
    title=#1,
    fonttitle=\bfseries,
    boxrule=0.8pt,
    arc=2mm
}
\newtcolorbox{ejemplo}[1]{
    colback=gray!8,
    colframe=gray!70!black,
    title=#1,
    fonttitle=\bfseries,
    boxrule=0.8pt,
    arc=2mm
}
\newtcolorbox{importante}[1]{
    colback=orange!8,
    colframe=orange!70!black,
    title=#1,
    fonttitle=\bfseries,
    boxrule=0.8pt,
    arc=2mm
}

\geometry{margin=1in}
\onehalfspacing
\pagestyle{fancy}
\rhead{Macroeconomía II}
\lhead{Clase 1}
\title{Clase 1}
\author{Ignacio Sanabria}
\date{Agosto 2026}

\begin{document}

\maketitle
\section{Conceptos}
\begin{itemize}
    \item Las economías suelen mantener una fluctuación entre la producción real, y la tendencia potencial. 
    \item Cuando el PIB real es menor que el potencial, se conoce como resecion. 
    
\end{itemize}
\section{Modelo Intertemporal de Consumo}
Suponga: 
\begin{itemize}
    \item Economía de dotación, no hay producción.
    \item Dos periodos $t\xrightarrow{}t+1$
    \item Agente representativo, un hogar ($H$) con una dotación \{$Y_t,Y_{t+1}$\}, en términos reales. 
    \item El numerario son las unidades de consumo. (unidad de medida).
    \item Queremos modelar el $C_t$ tal que: $$C_t=f(y_t,y_{t+1},...)$$
    \item Un mal modelo seria: $C_t=Y_t$ , dado que esta gastando todo su ingreso en el presente sin prever el futuro. 
    \item En $t$ el hogar decide $C_t$ , automáticamente decide $S_t$ (stock de ahorro) e intuitivamente $C_{t+1}$ 
\end{itemize}
Restricciones presupuestarias: 
\[\underbrace{C_t+S_t}_{Gastos} = \underbrace{Y_t}_{Ingresos} \]
\begin{itemize}
    \item Suponga que $S_{t-1}=0$
    \item Sea $r_t=$ la tasa de interés real.
    \item El hogar decide $S_t$ en el periodo $t$ y recibe $S_t+r_t\cdot S_t = (1+r_t)S_t $
    \item Si $S_t<0 \xrightarrow{}$ deuda. 
\end{itemize}
Por lo tanto, en el periodo $t+1$, la restricción presupuestaria es:
%FOrma de igualar los Iguales, en linea, una al lado del otro segun igual%
\begin{align*}
    C_t+S_t & =Y_t \\ C_{t+1}+S_{t+1} & =Y_{t+1}+(1+r_t)S_t
\end{align*}

Considere un hogar con preferencias: 
$$U=u(c_t)+\beta \cdot u(c_{t+1})$$
Tal que $u(\cdot)$ es una función con $$u`(\cdot ) >0\\u``(\cdot ) >0$$ 
Con $0<\beta<1$ , un factor de cambio. \\
Por lo tanto, podemos definir que cada agente busca: 
\begin{align*}
    \underbrace{max}_{\{ C_t, C_{t+1},S_t,S_{t+1}\}} u(C_t)+\beta\cdot u(C_{t+1}) \\ \text{s.a} \\ C_t+S_t & =Y_t \\  C_{t+1}+S_{t+1} & =Y_{t+1}+(1+r_t)S_t 
\end{align*}
Note que $S_{t+1}>0$ , no va a elegirlo pues tiene $u`(\cdot )>0$ , es decir no le genera utilidad no consumir todo. \\
Combinando las restricciones considerando esto, obtenemos: 
$$C_{t}+\frac{C_{t+1}}{1+r_t}=Y_t+\frac{Y_{t+1}}{1+r_t}$$
Donde $\frac{1}{1+r_t}$ es el factor de valor presente del modelo.\\  %Checar Notas% 
Ademas considere que el consumo son variables endógenas, y el ingreso presente y futuro junto con la tasa de interés son variables exógenas. \\

\subsection{Resolviendo:}
\[ 
\begin{aligned}
    \mathcal{L} &= u(c_t)+ \beta u(c_{t+1}) + \lambda [\;Y_t+\frac{Y_{t+1}}{1+r_t}-(C_{t}+\frac{C_{t+1}}{1+r_t})\; ] \\\\
    &\text{CPO:}\quad\\
    & \\  
    [c_t]:\quad &u'(c_t)-\lambda=0  \\
    [c_{t+1}]: \quad & u'\beta(c_{t+1})- \frac{\lambda }{1+r_t}=0\\
    [\lambda]:\quad &  Y_t+\frac{Y_{t+1}}{1+r_t}=C_{t}+\frac{C_{t+1}}{1+r_t}
\end{aligned}
\]
Por lo tanto igualando las COP podemos obtener la ecuación de Euler de forma que: 
\begin{align*}
    \frac{u'(c_t)}{u'(c_{t+1})}={\beta}({1+r_t})
\end{align*}

\begin{ejemplo}{Ejemplo 1}
    Suponga que $u(c)=log(c)$ tal que: $u'(c)=\frac{1}{c}$\\
    Las condiciones de optimalidad: 
    \begin{align}
        \frac{C_{t+1}}{1+r_t}=\beta(1+r_t)\\
        C_{t}+\frac{C_{t+1}}{1+r_t}=Y_t+\frac{Y_{t+1}}{1+r_t}
    \end{align}
    Al inyectar (1) en (2). \\
    $$ C_t=(\frac{1}{1+\beta})[Y_t+\frac{Y_{t+1}}{1+r_t}] $$
    El hogar consume una fraccion $\frac{1}{1+\beta}<1$ del valor del flujo de ingreso. Esto se conoce como \textbf{propensión marginal a consumir.} 
\end{ejemplo}


\subsection{Ecuación de Euler}
La ecuación de Euler básica es: 
$$u'(c_t)=\beta(1+r_t)+u'(c_{t+1})$$
Comportamiento de la ecuación de Euler: 
\begin{itemize}
    \item $\beta(1+r_t)=1$, podemos definir que: $u'(c_t)=u'(c_{t+1})$
    \item $\beta(1+r_t)<1$, podemos definir que: $u'(c_t)<u'(c_{t+1})\Rightarrow{} c_t>c_{t+1}$ \\\textit{En este caso la impaciencia $\beta$ domina la fuerza para postergar el consumo}
    \item $\beta(1+r_t)>1$, podemos definir que: $u'(c_t)>u'(c_{t+1})\Rightarrow{} c_t<c_{t+1}$
\end{itemize}

\subsection{Critica a Keynes}
\begin{comentarioprofe}{Critica}
    Según Keynes (1936) el consumo optimo se definía como $C_t=c_0+c_1Y_t$. \\
    Hecho estilizado (hasta 80s): $\frac{C_t}{Y_t} \sim 0.6-0.7$ \\
    Sin embargo, con preferencias logarítmicas sabemos que: 
    $$C_t=(\frac{1}{1+\beta})[Y_t+\frac{Y_{t+1}}{1+r_t}]$$
    Con esto yo puedo prever el accionar a futuro, por eventualidades o cambios abruptos en la tasa de interés ($r_t$)\\ 
    En este caso $PMC=\frac{1}{1+\beta}$, es decir se necesitan hogares muy impacientes para tener un mayor consumo. \\
    Si tomamos $\beta(1+r_t)=1$, sabemos que: 
    $$C_t=(\frac{1+r_t}{2+r_t})[Y_t+\frac{Y_{t+1}}{1+r_t}]$$ 
    donde el $PMC=\frac{1+r_t}{2+r_t}$
\end{comentarioprofe}

\begin{ejemplo}{Ejemplo 2}
    Suponga que posee un modelo $t+1$, tal que su linea de vida se ve: 
    \begin{figure} [H]
        \centering 
        \includesvg[width=0.7\linewidth]{figuras/clase_1/macro_1.drawio.svg}
        \label{fig:placeholder}
    \end{figure}
    Restricciones presupuestarias:
    \begin{align*}
        C_t+S_t=&Y_t\\
        C_{t+1}+S_{t+1}=&Y_{t+1}+(1+r_t)s_t\\ 
        C_{t+2}+S_{t+2}=&Y_{t+2}+(1+r_{t+2})S_{t+1}\\
        \vdots \\ 
        C_{t+T}+S_{t+T}=&Y_{t+T}+(1+r_{t+T-1})S_{t+T-1}\\\\
        \text{Condición terminal:}\; S_{t+T}&=0
    \end{align*}
    Para cualquier periodo $j$ tenemos: 
    \begin{align*}
        C_{t+j}+S_{t+j} =& Y_{t+j}+(1+r_{t+j-1})S_{t+j-1}\\
        C_{t+j}+\underbrace{S_{t+j}-S_{t+j-1}}_{\textbf{Flujo de ahorro}} =&Y_{t+j}+\underbrace{(r_{t+j-1})S_{t+j-1}}_{\textbf{Renta Ahorro}}
    \end{align*}
    El flujo de ahorro y la renta del ahorro a lo largo del tiempo se puede ver de la siguiente manera:
    \begin{figure} [H]
        \centering 
        \includesvg[width=0.8\linewidth]{figuras/clase_1/macro_2.drawio.svg}
        \label{fig:placeholder}
    \end{figure}
    \end{ejemplo}
    \newpage
    \begin{ejemplo}{Ejemplo 2}
        En el flujo del tiempo se pueden consolidar las restricciones presupuestarias y asumiendo que $r_{t+j}=r\; \forall\; j=0,1,...,T$ \\
        Podemos concretar que: 
        \begin{align*}
            C_t &=\frac{C_{t+1}}{1+r}+\frac{C_{t+2}}{(1+r)^2}+...+ \frac{C_{t+T}}{(1+r)^T}=Y_t+\frac{Y_{t+1}}{1+r}+...+\frac{Y_{t+T}}{(1+r)^T} \\\\
            &\text{Equivalente a la}\;\textbf{Restricción Presupuestaria Intertemporal (RPI):}\\\\
            &\sum^T_{j=0}\frac{C_{t+1}}{(1+r)^j}=\sum^T_{j=0}\frac{Y_{t+1}}{(1+r)^j}\\\\
            &\textbf{Riqueza o Ingreso Permanente:}\; =W=\sum^T_{j=0}\frac{Y_{t+1}}{(1+r)^j}
        \end{align*}
        Si intentamos maximizar el problema de utilidad para $T$ periodos obtenemos: 
        \begin{align*}
            \underbrace{max}_{ \{ C_{t+j},Y_j\}\ }  u(C_t) + \beta \;u(C_{t+1})+\beta^2 \;u(C_{t+2})+...+\beta^T \;u(C_{t+T})\\
            \text{Esto es equivalente a:}\\ \sum_{j=0}^T\beta^j u(C_{t+j})\quad \text{s.a}\quad \sum^T_{j=0}\frac{C_{t+1}}{(1+r)^j}=\sum^T_{j=0}\frac{Y_{t+1}}{(1+r)^j}\\\\
            \text{En un lagrangiano se debería derivar con respecto a cada $\beta$}\\ \text{En este caso existirían $T$ Ecuaciones de Euler}
        \end{align*}
    \end{ejemplo}
\newpage
Ahora suponga que $\beta(1+r_t)=1$, en este caso como vimos con anterioridad el consumo sera constante durante todo su ciclo de vida. Por lo cual podríamos definir un $\bar {C}$. 
\begin{align*}
    \text{Suponga que:}\; \beta(1+r_t)=1 \\
    \text{Por lo tanto definimos:} \\
    C_t=C_{t+1}=&C_{t+2}=...=C_{t+T}=\mathbf{\bar{C}}\\\\
    \text{Ahora cambiemos esto en nuestra }\textbf{RPI}\\\\
    \sum^T_{j=0}\frac{C_{t+1}}{(1+r)^j}=\sum^T_{j=0}\frac{Y_{t+1}}{(1+r)^j} & \Rightarrow{} \sum^T_{j=0}\frac{\mathbf{\bar{C}}}{(1+r)^j}=\sum^T_{j=0}\frac{Y_{t+1}}{(1+r)^j} \\\\
    \mathbf{\bar{C}} \cdot \sum^T_{j=0}\frac{1}{(1+r)^j}=\sum^T_{j=0}\frac{Y_{t+1}}{(1+r)^j} & \Rightarrow{} \mathbf{\bar{C}} \cdot \sum^T_{j=0} \beta^j=\sum_{j=0}^T \beta^j Y_{t+j}\\\Rightarrow{} \mathbf{\bar{C}} \cdot (\frac{1-\beta^{T+1}}{1-\beta})=\sum_{j=0}^T \beta^j Y_{t+j} \\\\ 
    \therefore \mathbf{\bar{C}} = (\frac{1-\beta^{T+1}}{1-\beta}) \sum_{j=0}^T \beta^j Y_{t+j}\equiv & (\frac{1-\beta^{T+1}}{1-\beta}) \cdot W
\end{align*}
De aqui podemos discernir también la PMC del individuo a lo largo del tiempo, ya que observamos que: 
$$\frac{\delta C_t}{\delta Y_t}=\frac{\delta \bar{C}}{\delta Y_t}=\frac{1-\beta^{T+1}}{1-\beta}\; \in (0,1)$$
\begin{importante}{Concluimos}
    \begin{itemize}
        \item Los hogares toman decisiones de consumo en función de su ingreso permanente, \textit{Hipótesis de Renta Permanente de Friedman}. 
        \item La PMC es \textbf{menor} entre mayor sea $T$
        \item $\frac{\delta C_t}{\delta Y_{t+1}}=\frac{\delta \bar{C}}{\delta Y_{t+1}}=\beta^j(\frac{1-\beta}{1-\beta^{t+1}})$ \\ 
        Choque de ingreso futuro mueve el consumo presente. pero entre mas $\underbrace{\text{alejado del presente}}_{j\;\text{mas alto}},\quad  \underbrace{\text{menos lo mueve}}_{\frac{\delta \bar{C}}{\delta Y_{t+1}}\;\text{es menor}}.$\\
    \end{itemize}
\end{importante}
\subsection{Tipos de Choque}
Note que: $d\bar{C}=(\frac{1-\beta^{T+1}}{1-\beta})(dY_t+\beta \;dY_{t+1}+\beta ^2 \; dY_{t+2}+...+\beta ^T\; dY_{t+T})$\\
De aquí podemos partir a clasificar los choques, por ejemplo> 
\subsubsection{Choque Completamente Transitorio}
$dY_t \neq 0, \; dY_{t+j}=0 \quad \forall j=1,2,...,T$\\
Implicando entonces que: $$d\bar{C}=(\frac{1-\beta^{T+1}}{1-\beta}) dY_t$$\\ 
es decir $$\frac{d\bar{C}}{dY_t}=(\frac{1-\beta^{T+1}}{1-\beta})=\textbf{PMC}$$
\subsubsection{Choque Completamente Permanente}
$dY_t \neq 0, \; dY_t=dY_{t+1}= dY_{t+2}=...= dY_{t+T}$\\
Implicando entonces que: 
\begin{align*}
    d\bar{C}=(\frac{1-\beta^{T+1}}{1-\beta}) (dY_t+\beta \;dY_t+\beta ^2 \; dY_t+...+\beta ^T\; dY_t)\\
    d\bar{C}=(\frac{1-\beta^{T+1}}{1-\beta})\;dY_T (1+\beta+\beta^2+...+\beta^T)\\
    d\bar{C}=(\frac{1-\beta^{T+1}}{1-\beta})\;dY_T \sum_{j=0}^T\beta^j\Rightarrow{}d\bar{C}=(\frac{1-\beta^{T+1}}{1-\beta})\;dY_T\frac{1-\beta}{1-\beta^{T+1}}\\ 
    d\bar{C}=dY_t\\\\
    \therefore \frac{ d\bar{C}}{dY_t}=1 =\textbf{PMC}
\end{align*}
De aqui podemos concluir que: \textbf{entre más persistente el choque de ingreso, más fuerte es la respuesta del consumo}.
\newpage
\subsection{Dos periodos}
Sin ahorro las condiciones de optimalidad: 
\begin{itemize}
    \item $u'(C_t)=\beta (1+r_t)\; u'(C_{t+1})$
    \item $C_{t}+\frac{C_{t+1}}{1+r_t}=Y_t+\frac{Y_{t+1}}{1+r_t}$
\end{itemize}
Mientras que con ahorro las condiciones de otimalidad: 
\begin{itemize}
    \item $u'(C_t)=\beta (1+r_t)\; u'(C_{t+1})$
    \item $C_t+S_t=Y_t$ siempre que $S_{t-1}=0$
    \item $C_{t+1}=Y_{t+1}+(1+r_t)S_t$
\end{itemize}
Por ende podemos decir que existe una función de consumo: 
$$C_t=C^d(\underbrace{Y_t}_{(+)},\underbrace{Y_{t+1}}_{(+)},{r_t} )$$
\subsection{Choques de Tasa de Interés}
$C_{t+1}=(Y_t-C_t)(1+r_t)+Y_{t+1}$
Cuya pendiente es: $-(1+r_t)$
Y se ve de forma: 
\begin{figure} [H]
        \centering 
        \includesvg[width=0.5\linewidth]{figuras/clase_1/macro_3.drawio.svg}
        \label{fig:placeholder}
\end{figure}
De esta figura podemos apreciar las curvas de indiferencia segun el nivel de ahorro y el punto de Autarquia. \\
\textit{En autarquia se parte el grafico entre las partes que puede consumir sin ser deudor ni ahorrante}\\\\
Si tenemos un cambio en la tasa de interes de $r_1>r_0$ sabiendo que la pendiente de la curva (en rojo) es $-(1+r_0)$ y su nueva pendiente la curva (en azul) es $-(1+r_1)$
\begin{figure} [H]
        \centering 
        \includesvg[width=0.5\linewidth]{figuras/clase_1/macro_4.drawio.svg}
        \label{fig:placeholder}
\end{figure}
Podemos observar que al haber una mayor pendiente en $r_1$, este hogar representativo ve limitado su consumo hoy $(C_t)$ pero gana mayor espacio de consumo mañana $(C_{t+1})$. 
\begin{ejemplo}{Logarítmica}
    Si $u(c)=log(c)$, sabemos que: 
    $$C_t=(\frac{1}{1+\beta})[Y_t+\frac{Y_{t+1}}{1+r_t}]$$
    y que $Y_t=S_t+C_t$ \\
    Por lo tanto llegamos a que: 
    $$S_t=(\frac{1}{1+\beta})(\beta Y_t-\frac{Y_{t+1}}{1+r_t})$$
    Viendo esto sabemos que si $$S_t<0 \iff Y_t<\frac{Y_{t+1}}{1+r_t})$$
\end{ejemplo}

Basándonos en la RPI, la relación de $C_t, r_t$ es ambiguo. 
\begin{itemize}
    \item Si es ahorrante, $\uparrow r_t, \uparrow C_t$
    \item Si es deudora, $\uparrow r_t, \downarrow C_t,\downarrow C_{}$
    \item Esto es la representación del \textbf{efecto ingreso}. 
\end{itemize}
De la Ecuacion de Euler sabemos: 
$u'(C_t)=\beta (1+r_t)\; u'(C_{t+1})$
Por ende: 
\begin{itemize}
    \item Si $\uparrow r_t\Rightarrow{} \downarrow C_t, \uparrow C_{t+1}$
    \item Sustituye el bien caro por un bien mas barato. 
    \item Esto se conoce como el \textbf{efecto sustitucion}.
\end{itemize}
En conclusion: \textit{El efecto de $r_t$ sobre $c_t$ depende de si el ES domina el EI}
$$C_t=(\underbrace{Y_t}_{(+)},\underbrace{Y_{t+1}}_{(+)},\underbrace{r_t}_{(?)} )$$

\begin{ejemplo}{Logarítmica}
    Si $u(c)=log(c)$, sabemos que: 
    $$C_t=(\frac{1}{1+\beta})[Y_t+\frac{Y_{t+1}}{1+r_t}]$$
    En este caso tenemos que:
    $$\frac{\partial C_t}{\partial r_t}<0 \xrightarrow{} \text{Efecto Sustitución}$$
\end{ejemplo}
\newpage
\begin{ejemplo}{Funcion CARRA}
    Sea: $$u(c)=\frac{c^{1-\frac{1}{\sigma}}}{1-\frac{1}{\sigma}},\;\; \sigma>0$$
    En este caso: $u'(c)=c^{\frac{-1}{\sigma}}$\\
    Tomemos en consideracion que si $\sigma\xrightarrow{}1$, se convierte en el ejemplo logaritmico ya visto.\\
    Por ende sabemos que la \textbf{Ecuacion de Euler} esta dada por: \\
    $\frac{C_{t+1}}{C_t}=[\beta(1+r_t)]^\sigma$\\
    \begin{itemize}
        \item Si $\beta(1+r_t)=1\; \Rightarrow{} \frac{C_{t+1}}{C_t}=1 \Rightarrow{} C_{t+1}=C_t$
        \item Si $\beta(1+r_t)>1\; \Rightarrow{} \frac{C_{t+1}}{C_t}>1$
    \end{itemize}
    Por ende sabemos que $\frac{C_{t+1}}{C_t}=[\beta(1+r_t)]^\sigma$\\
\end{ejemplo}
Recordemos la definicion de una funcion de cambio porcentual, o crecimiento. 
\begin{align*}
    \text{Sea:}\; X_t,X_{t-1} \\
    g_t^x= \frac{X_t-X_{t-1}}{X_{t-1}} \\ 
    g_t^x= \frac{X_t}{X_{t-1}}-1 \\
    1-g_t^x=\frac{X_t}{X_{t-1}}\\\\
    \text{Sea $t=t+1$};\\
    1-g_{t+1}^x=\frac{X_{t+1}}{X_{t}}\\\\
\end{align*}
\begin{ejemplo}{Función CARRA}
    Ahora que definimos una función de crecimiento podemos comparar:
       $$ \frac{X_{t+1}}{X_{t}}=1-g_{t+1}^x=\frac{C_{t+1}}{C_t}$$
    Por ende definimos:
    $$\frac{C_{t+1}}{C_t}=1-g_{t+1}^x$$
    Siendo $g_{t+1}^x$ la tasa de crecimiento de consumo futuro. \\\\
    Ademas si agregamos lgaritmo a ambos lados de la ecuacion:\\
    \begin{align*}
        log(\frac{C_{t+1}}{C_t})=log([\beta(1+r_t)]^\sigma)\\
        log(\frac{C_{t+1}}{C_t})=\sigma log(\beta)+\sigma log(1+r_t)\\
        \sigma=\frac{\partial log(\frac{C_{t+1}}{C_t})}{\partial log(1+r_t)}
    \end{align*}
\end{ejemplo}
En este caso $\sigma$ es la \textbf{elasticidad sustitucion intertemporal}, es decir, que tanto cambia $\frac{C_{t+1}}{C_t}$ cuando aumenta $(1+r_t)$\\

\begin{importante}{Observaciones}
    \begin{itemize}
    \item $\sigma$ es la elasticidad entre como cambia el consumo del hogar presente y futuro, según el cambio en el interés. 
    \item Si $\sigma$ es muy alto el cambio en el consumo del hogar es muy fuerte. 
    \item $\sigma$ es una elasticidad de sustitución intertemporal, que tan fuerte es el efecto sustitución. 
    \item Empíricamente el $\sigma$ no afecta tanto en el consumo de los hogares, es muy baja la relación entre el comportamiento de la tasa de interés real y el consumo del hogar. 
\end{itemize}
\end{importante} 
\begin{ejemplo}{Funcion CARRA}
    Sea $$u(c)=\frac{c^{1-\frac{1}{\sigma}}}{1-\frac{1}{\sigma}},$$. Y su derivada sea $u'(c)=c^{-\frac{1}{\sigma}}$\\
    Aplicando el lagrangiano con la RPI obtenemos:
    \begin{align*}
        [C_t]: C_t^{-\frac{1}{\sigma}}-\lambda=0\\
        [C_{t+1}]: C_{t+1}^{-\frac{1}{\sigma}}\beta-\frac{\lambda}{1+r_t}=0\\
        {(\frac{C_{t+1}}{C_{t}}})^{-\frac{1}{\sigma}}=\frac{\frac{\lambda}{\beta(1+r_t)}}{\lambda}\\
        {(\frac{C_{t+1}}{C_t}})^{-{1}}=\frac{1}{(\beta(1+r_t))^\sigma}\\
        \frac{C_{t+1}}{C_t}={(\beta(1+r_t))^\sigma}\\
        C_{t+1}={(\beta(1+r_t))^\sigma} \cdot C_t
    \end{align*}
    Ahora procedemos a inyectar $C_t$ en la RPI
    \begin{align*}
        C_t+\frac{{(\beta(1+r_t))^\sigma} \cdot C_t}{(1+r_t)}=(Y_t+\frac{Y_{t+1}}{1+r_t})\\
        C_t(1+\frac{{(\beta(1+r_t))^\sigma}}{(1+r_t)})=(Y_t+\frac{Y_{t+1}}{1+r_t})\\
        C_t(\frac{(1+r_t)+(\beta(1+r_t))^\sigma}{(1+r_t)})=(Y_t+\frac{Y_{t+1}}{1+r_t})\\
        C_t=[\frac{(1+r_t)}{(1+r_t)+(\beta(1+r_t))^\sigma}](Y_t+\frac{Y_{t+1}}{1+r_t})
    \end{align*}
    La forma de comprobar que esta igualdad es correcta es si $\sigma=1$ obtenemos el ejemplo Logaritmica donde $\frac{\partial C_t}{\partial Y_t}=(\frac{1}{1+\beta})$
\end{ejemplo}
\begin{ejemplo}{Funcion CARRA}
    Podemos comprobar que: 
    \begin{align*}
        \text{Si $\sigma =1$ obtenemos:}\\
        C_t=[\frac{(1+r_t)}{(1+r_t)+(\beta(1+r_t))^1}](Y_t+\frac{Y_{t+1}}{1+r_t})\\
        C_t=\frac{(1+r_t)}{(1+r_t)(1+(\beta))}(Y_t+\frac{Y_{t+1}}{1+r_t})\\
        \therefore C_t=(\frac{1}{1+\beta})(Y_t+\frac{Y_{t+1}}{1+r_t})\\\\
        \frac{\partial C_t}{\partial Y_t}=(\frac{1}{1+\beta})
    \end{align*}
    De aqui podemos ver que el resultado de $\frac{\partial C_t}{\partial r_t}$ depende de $\sigma$. \\
    Tal que: 
    \begin{equation*}
        \begin{array}{cc}
         \text{Si}\; \sigma \le 1 & \Rightarrow{}   \frac{\partial C_t}{\partial r_t} <0\\
         \text{Si}\; \sigma>1 & \Rightarrow{}   \frac{\partial C_t}{\partial r_t} >0
    \end{array}
    \end{equation*}
    
\end{ejemplo}

\newpage
\section{Agregación y Equilibrio General}
Suponga dos tipos de hogares con preferencias distintas, cada hogar tiene medida $N1$ y $N2$. 
Entonces: 
\begin{align*}
    C_t=N1\;C_t^1+N2\;C_t^2\\
    S_t=N1\;S_t^1+N2\;S_t^2
\end{align*}
Para comprender mejor veamos el siguiente ejemplo: 
\begin{ejemplo}{2 Hogares Representativos}
    Suponga: 
    \begin{align*}
        (Y_t^1,Y_{t+1}^1)=(1,0)\\
        (Y_t^2,Y_{t+1}^2)=(0,1)\\
    \end{align*}
    Suponga ahora que ambos hogares tienen las mismas preferencias, tal que: 
    \begin{align*}
        u(C_{t+1}^1,C_t^1)=u(C_t^1)+\beta u(C_{t+1}^1)\\
        u(C_{t+1}^2,C_t^2)=u(C_t^2)+\beta u(C_{t+1}^2)\\
    \end{align*}
    Cuyas funciones de consumo dado su ingreso son: 
    \begin{align*}
        C_t^1=(\frac{1}{1+\beta})(Y_t^1+\frac{Y_{t+1}^1}{1+r_t})&=(\frac{1}{1+\beta})\\
        C_t^2=(\frac{1}{1+\beta})(Y_t^2+\frac{Y_{t+1}^2}{1+r_t})& =(\frac{1}{1+\beta})(\frac{1}{1+r_t})\\
    \end{align*}
    El ahorro de estos hogares se veria: 
    \begin{align*}
        S_t^1=&(\frac{\beta}{1+\beta})(Y_t^1-\frac{Y_{t+1}^1}{1+r_t})=\frac{\beta}{1+\beta}\\
        S_t^2=&(\frac{\beta}{1+\beta})(Y_t^2-\frac{Y_{t+1}^2}{1+r_t})=(\frac{\beta}{1+\beta})(\frac{1}{1+r_t})
    \end{align*}
\end{ejemplo}
\begin{ejemplo}{2 Hogares Representativos}
    Suponga entonces que N1=1, N2=1\\
    $$S_t(r_t)=\frac{\beta}{1+\beta}-(\frac{\beta}{1+\beta})(\frac{1}{1+r_t})$$
    Por ende: \\
    \begin{equation*}
        \begin{array}{cc}
            \text{Si } S_t(r_t)>0 & \Rightarrow{} \text{Exeso de oferta de fondos prestables}  \\
            \text{Si } S_t(r_t)<0  & \Rightarrow{} \text{Exeso de demanda de fondos prestables} 
        \end{array}
    \end{equation*}
    Entonces: $r_t^*$ vacia o aclara el mercado de fondos prestables si no hay exceso de oferta ni demanda. \\
    Generando que $$S_t(r_t^*)=0$$
\end{ejemplo}

\end{document}